\documentclass[11pt]{article}
\usepackage[letterpaper,margin=1in]{geometry}
\usepackage[T1]{fontenc}
\usepackage{lmodern}
\usepackage{amsmath,amssymb}
\usepackage[colorlinks=true,linkcolor=blue,urlcolor=blue,citecolor=blue]{hyperref}
\usepackage{microtype}

\title{A Contractively Complemented Approximate Identity:\\
Literature Resolution of arXiv:math/0512417}
\author{Run \texttt{fa\_banach\_001} --- literature-status note}
\date{August 11, 2026}

\begin{document}
\maketitle

\noindent\textbf{Status.} \texttt{literature\_already\_answered}. This note
records an explicit later resolution; it does not claim a new theorem.

\section*{Source question}
Blecher, Hay, and Neal ask in Section 6 of \cite{BHN} (arXiv PDF page 16;
published page 351):
\begin{quote}
Does every approximately unital operator algebra $A$ have an approximate
identity of the form $(1-x_t)$ with $x_t\in\operatorname{Ball}(A^1)$?
\end{quote}
Here $A^1$ is the unitization. They explain that this is equivalent to the
question whether every approximate $p$-projection is a $p$-projection.

\section*{Decisive later answer}
Blecher and Read explicitly state that their Theorem 1.1 was ``prompted by,
and solves'' the question on page 351 of \cite{BHN}. Their theorem says that
every operator algebra with a contractive approximate identity has one,
say $(e_t)$, satisfying
\[
 \|1-e_t\|\le 1,
 \qquad\text{and even}\qquad
 \|1-2e_t\|\le 1
 \quad\text{for every }t \cite[Theorem 1.1]{BR}.
\]
Set $x_t=1-e_t\in A^1$. Then $e_t=1-x_t$ and
$\|x_t\|=\|1-e_t\|\le1$, exactly as requested. The stronger second
inequality also implies the first bound by
\[
 x_t=\frac{1+(1-2e_t)}2.
\]
The same paper states the equivalent geometric conclusion: every closed
projection in $A^{**}$ is a $p$-projection and is a strong limit of a
decreasing net of peak projections \cite[Theorem 1.2]{BR}.

\section*{Provenance and scope}
This is an exact, author-recognized answer rather than an implication first
noticed here: the later paper cites the original source and identifies its
page number. It settles both the displayed approximate-identity question and
the equivalent $p$-projection question discussed with it. No residual part of
that extracted question remains open. Other broad questions in the source
paper, such as general proximinality of right ideals, are separate and are not
claimed to be settled by this identification.

\begin{thebibliography}{9}
\bibitem{BHN}
D. P. Blecher, D. M. Hay, and M. Neal,
\emph{Hereditary subalgebras of operator algebras},
J. Operator Theory 59 (2008), 333--357; arXiv:math/0512417.

\bibitem{BR}
D. P. Blecher and C. J. Read,
\emph{Operator algebras with contractive approximate identities},
arXiv:1011.3797 (2010, revised 2011).
\end{thebibliography}

\end{document}

